% =====================================================================
% Apêndice A: código-fonte (lido direto de codigo/rlc_runge_kutta.py)
% =====================================================================
\appendix
\section{Código-fonte da implementação}
\label{ap:codigo}

O código abaixo é o programa executado para produzir as tabelas e as figuras do trabalho. O espaçamento do apêndice é simples para que a listagem permaneça legível; o corpo do trabalho usa espaçamento 1,5. A linguagem é Python 3. Os métodos de Euler, de Euler melhorado, de Runge-Kutta de quarta ordem e o controlador de passo estão implementados nas funções \texttt{passo\_euler}, \texttt{passo\_heun}, \texttt{passo\_rk4} e \texttt{rk4\_adaptativo}.

\begin{singlespace}
\VerbatimInput[%
  fontsize=\footnotesize,%
  numbers=left,%
  numbersep=6pt,%
  breaklines=true,%
  breakanywhere=true%
]{codigo/rlc_runge_kutta.py}
\end{singlespace}
